\documentclass[11pt]{article}
\usepackage[T1]{fontenc}
\usepackage{lmodern,amsmath,amssymb,bm,geometry,booktabs,longtable,hyperref}
\geometry{margin=24mm}
\hypersetup{colorlinks=true,linkcolor=blue,urlcolor=blue}
\setlength{\emergencystretch}{3em}
\newcommand{\NS}{\mathrm{NS}}
\newcommand{\R}{\mathrm R}
\newcommand{\F}{\mathsf F}
\numberwithin{equation}{section}
\title{Superconformal blocks on a plumbing graph}
\author{Supplementary implementation notes}
\date{September 16, 2026}
\begin{document}
\maketitle

A superconformal block is a sum of three-point forms sewn with inverse
Gram matrices. We compute it by adjoining a free fermion, reducing the
enlarged block to a sum of products of two Virasoro blocks, and removing
the fermion factor by convolution. The challenge is to keep the operator order and local
normalizations consistent through those three steps.

We first specify the PBW sewing. We then derive the local fermion signs,
separate them into physical, auxiliary, and mixed factors, and obtain the
convolution. The double-Virasoro formula evaluates its enlarged side.
Finally, the Ramond cycle identities identify the sector on which recovery
is possible; one external $\Theta\psi$ insertion supplies a missing sector.
For Mercedes with NS spokes and a Ramond rim, this requires at most one
inserted sphere, while the central all-NS sphere stays sewn in place.

We keep the notation of the SCblock notes and the draft. In particular,
$K$ denotes the operator-reordering exponent and $\star$ its associated
convolution. The formulas below display the state arguments and local
signs explicitly; no separate state aliases or parity-transformation
operators are needed.

\section{States, three-point forms, and physical sewing}

\subsection{The channel and its labels}

A plumbing channel $\hat\Gamma$ specifies the NS/R edges, local coordinates,
and spin lifts $\eta_e=\pm1$. Each trivalent sphere is either all-NS or
NS--R--R. Fix the slots $(\infty,1,0)$, numbered $1,2,3$; at an NS--R--R
sphere slot 1 is NS. Also fix an order of the spheres, the internal edges,
the two ends of each edge, and the external punctures. These choices
specify the operator order as well as the sewing. Changing a slot order
requires transporting its states, coordinates, and plumbing maps together.
A self-loop has two distinct ends at the same vertex.

Subscripts $(e,1),(e,2)$ label the ends of an internal edge, while
$(v,1),(v,2),(v,3)$ label the slots of a sphere. They refer to the same
punctures: for example, $A_{v,i}$ is the PBW word at the end of the edge
incident to that slot. Throughout, edge products and internal-state sums
run only over internal edges unless stated otherwise. External states are
fixed. Products restricted by $e:\NS$, $e:\R$, or
$v:\NS,\R,\R$ specify the corresponding edge or vertex types, without
introducing separate sets for each type.

For a connected trivalent channel the genus is the number of internal
edges minus the number of vertices plus one. The elementary unstable
sphere pairings are treated separately; an unpunctured torus can first
be given an identity puncture. Equivalent local spin frames do not count
as additional spin structures.

\subsection{The PBW basis and local forms}

An NS PBW state is $\mathbf L_{-A}\phi$, and a Ramond state is
$\mathbb L_{-A}w^\alpha$, where the word contains ordered $L_{-n}$ and
$G_{-r}$ with positive mode indices. The $G$ indices are half-integer in
NS and integer in R. As in the draft, $|A|$ is the sum of the mode indices,
whereas $A$ in a sign counts the $G$ operators modulo two. The NS primary
parity is $p$; the two physical Ramond grounds are $w^+$ even and $w^-$
odd. Thus
\begin{equation}
 \epsilon=\begin{cases}A+p,&\NS,\\A+|\alpha|,&\R,\end{cases}
 \qquad |+|=0,\quad|-|=1.
 \label{eq:physical-states}
\end{equation}
Every parity is represented by zero or one. Parity sums are modulo two
unless a different meaning is explicitly specified.
The Ramond parameter obeys $\beta^2=c/24-h$, where $c$ is the physical
superconformal central charge, and
$G_0w^\pm=i\beta e^{\mp i\pi/4}w^\mp$.
Normalize the NS primary to BPZ norm one and use the physical Ramond
ground metric $\operatorname{diag}(1,i)$.

Lower Gram indices denote pairings and upper indices their inverses:
\begin{align}
 (\mathbf G_h)_{AB}
 &=\langle\!\langle\mathbf L_{-A}\phi_h|
                                     \mathbf L_{-B}\phi_h\rangle,
 \nonumber\\
 (\mathbb G_\beta)_{A\alpha,B\gamma}
 &=\langle\!\langle\mathbb L_{-A}w_\beta^\alpha|
                                     \mathbb L_{-B}w_\beta^\gamma\rangle.
 \label{eq:physical-metric}
\end{align}
They pair equal levels and equal parities, so a nonzero edge contraction
has one common level and one common $\epsilon_e$ at its two ends.
These are bilinear BPZ pairings, without complex conjugating coefficients.
We assume generic weights for which the inverse Gram matrices exist.
Degenerate modules require their null-state quotient or a specified
regular limit.

At a vertex the label is $a_v=0,1$ for all-NS, or
$f_v=0,1$ and $\eta_v=\pm1$ for NS--R--R. The vertex sign $\eta_v$
labels a three-point form; the edge sign $\eta_e$ labels sewing.
The forms at $(\infty,1,0)$ are normalized by
\begin{align}
 \rho_0(\phi_1,\phi_2,\phi_3)
 &=\rho_1(\phi_1,G_{-1/2}\phi_2,\phi_3)=1,\nonumber\\
 \rho_0^{(\eta)}(\phi_1,w_2^+,w_3^+)
 &=\rho_1^{(\eta)}(\phi_1,w_2^+,w_3^-)=1,\nonumber\\
 \rho_0^{(\eta)}(\phi_1,w_2^-,w_3^-)&=\eta,
 \qquad\rho_1^{(\eta)}(\phi_1,w_2^-,w_3^+)=i\eta.
 \label{eq:physical-vertices}
\end{align}
Unlisted Ramond ground components vanish. Superconformal Ward identities
in this slot order determine their descendants. Nonzero components obey
\begin{equation}
 \sum_i\epsilon_{v,i}
 =\begin{cases}a_v+\sum_i p_{v,i},&(\NS,\NS,\NS),\\
                f_v+p_{v,1},&(\NS,\R,\R).
   \end{cases}
 \label{eq:vertex-selection}
\end{equation}
Equivalently, the descendant and ground labels obey $A+B+C=a$ in all-NS
and $A+B+C+|\alpha|+|\gamma|=f$ in NS--R--R. These constraints are
used below to simplify the signs.

\subsection{The operator-reordering sign and the block}

PBW sewing initially gives consecutive pairs of operators, one pair per
internal edge, followed by the external operators. To evaluate the local
three-point forms, reorder them into consecutive vertex triples.
The exponent $K$ records the odd exchanges:
\begin{equation}
 K(\epsilon)=\sum_{\substack{
 (v,i)\text{ precedes }(w,j)\text{ in edge-pair order}\\
 (w,j)\text{ precedes }(v,i)\text{ in vertex order}}}
                         \epsilon_{v,i}\epsilon_{w,j}.
 \label{eq:graph-sign}
\end{equation}
The external punctures are included in the first order after all internal
pairs. Each reversed pair is counted once. For example, in theta the
orders $(1_L,1_R,2_L,2_R,3_L,3_R)$ and
$(1_L,2_L,3_L,1_R,2_R,3_R)$ give
$K=\epsilon_1\epsilon_2+\epsilon_1\epsilon_3+\epsilon_2\epsilon_3$.
The same prescription includes self-loops and external odd states without
referring to crossings in a drawing.

The physical block now follows by inserting each inverse Gram matrix
between its two vertex forms:
\begin{align}
 \mathbf F_{\hat\Gamma}={}&\sum_{\{A,\alpha\}}(-1)^{K(\epsilon)}
 \prod_{e:\NS}q_e^{|A_{e,1}|}\eta_e^{\epsilon_e}
            \mathbf G_{h_e}^{A_{e,1}A_{e,2}}
 \prod_{e:\R}q_e^{|A_{e,1}|}\eta_e^{\epsilon_e}
            \mathbb G_{\beta_e}^{A_{e,1}\alpha_{e,1},A_{e,2}\alpha_{e,2}}
 \nonumber\\
 &\times\prod_{v:\NS,\NS,\NS}
 \rho_{a_v}(\mathbf L_{-A_{v,1}}\phi_{v,1},
            \mathbf L_{-A_{v,2}}\phi_{v,2},
            \mathbf L_{-A_{v,3}}\phi_{v,3})\nonumber\\
 &\times\prod_{v:\NS,\R,\R}
 \rho_{f_v}^{(\eta_v)}(\mathbf L_{-A_{v,1}}\phi_{v,1},
            \mathbb L_{-A_{v,2}}w_{v,2}^{\alpha_{v,2}},
            \mathbb L_{-A_{v,3}}w_{v,3}^{\alpha_{v,3}}).
 \label{eq:physical-block}
\end{align}
The sum includes both physical endpoint words on every internal edge and
both Ramond ground labels where applicable. Equal-level and equal-parity
constraints are enforced by the inverse Gram entries. The local
constraints are enforced by the vertex forms. The powers $q_e^{h_e}$
are excluded from this descendant-normalized block and restored in a
correlator. External finite linear combinations are treated by linearity.

\section{The local signs and the enlarged block}
\label{sec:local-signs}

\subsection{The auxiliary basis and its BPZ pairing}

Adjoin an auxiliary fermion, with
$\{\psi_r,\psi_s\}=\delta_{r+s,0}$ and
$\psi_r^\dagger=-\psi_{-r}$. The minus sign in this BPZ conjugation is
required for the two Virasoro generators to obey
$(L_n^{(j)})^\dagger=L_{-n}^{(j)}$. Write
$\Psi_{-\mathsf A}=\psi_{-r_1}\cdots\psi_{-r_m}$ with
$r_1>\cdots>r_m>0$. The modes are half-integer in NS and integer in R;
$\mathsf A$ in a sign counts these modes, and $|\mathsf A|$ is their
sum. Zero modes are represented by the ground labels, not included in
$\Psi_{-\mathsf A}$. The auxiliary states and parities are
\begin{equation}
 \begin{array}{c|c|c}
 &\text{state}&\epsilon'\\ \hline
 \NS&\Psi_{-\mathsf A}\mathbf1&\mathsf A\bmod2\\
 \R&\Psi_{-\mathsf A}u^{\mathsf a}&(\mathsf A+\mathsf a)\bmod2
 \end{array}
 \qquad \mathsf a\in\{0,1\}.
 \label{eq:aux-states}
\end{equation}
Here $\epsilon$ continues to denote physical parity, and $\epsilon'$
denotes auxiliary parity; the tensor state has parity $\epsilon+\epsilon'$.
Choose $\psi_0u^{\mathsf a}=u^{1-\mathsf a}/\sqrt2$ and ground norms
$1,-1$. The oscillator anticommutators then give
\begin{equation}
 \langle\!\langle\Psi_{-\mathsf A}u^{\mathsf a}|
                         \Psi_{-\mathsf B}u^{\mathsf b}\rangle
 =\begin{cases}(-1)^{\mathsf A+\mathsf a},
                    &(\mathsf A,\mathsf a)=(\mathsf B,\mathsf b),\\
                0,&\text{otherwise}.
   \end{cases}
 \label{eq:aux-metric}
\end{equation}
For NS replace the grounds by $\mathbf1$ and set $\mathsf a=0$.
Thus the inverse auxiliary Gram entry is also $(-1)^{\epsilon'}$.
On the tensor basis we use the product of the physical and auxiliary BPZ
pairings, without an additional crossing sign. This local convention is
separate from the global operator permutation counted by $K$.

The normalized auxiliary three-point form is fixed by
\begin{equation}
 \rho_{\F}(\mathbf1,\mathbf1,\mathbf1)
 =\rho_{\F}(\mathbf1,u^0,u^0)=1,
 \qquad \rho_{\F}(\mathbf1,u^1,u^1)=i,
 \label{eq:aux-ground}
\end{equation}
and the fermion Ward identities. It vanishes unless the three auxiliary
parities sum to zero. The BPZ signs of bra descendants are already included
in $\rho_{\F}$; they must not be added a second time.

\subsection{Why the local three-point forms carry signs}

There are two distinct operations: exchanging odd operators, and converting
between the normalization of the physical forms and that of the
Virasoro branching states. An ordinary graded rearrangement from three
auxiliary--physical pairs to three auxiliary states followed by three
physical states would give
\begin{equation}
 (-1)^{\epsilon_1(\epsilon'_2+\epsilon'_3)+\epsilon_2\epsilon'_3}.
 \label{eq:literal-moving}
\end{equation}
The first physical state passes the second and third auxiliary states;
the second physical state passes the third auxiliary state. This accounts
for each term in the exponent. The enlarged forms below also incorporate
the fixed BPZ and contour conventions used for double Virasoro. They are
therefore not defined by this permutation sign alone.

For even NS primaries, \eqref{eq:literal-moving} becomes
$(-1)^{A\mathsf A+B\mathsf C}$. The product BPZ pairing removes the
crossing of the two factors in the bra, supplying a second
$(-1)^{A\mathsf A}$. Their product leaves $(-1)^{B\mathsf C}$.
It is the crossing of the middle-slot supercurrent modes with the ket
fermions, expressed directly by
\begin{equation}
 \mathbf L_{-B}\Psi_{-\mathsf C}
       =(-1)^{B\mathsf C}\Psi_{-\mathsf C}\mathbf L_{-B}.
 \label{eq:middle-ket-crossing}
\end{equation}
The intrinsic primary parities require
one further conversion. At fixed weights and descendant words, the
physical Ward convention obeys
\begin{align}
 &\rho_a(\mathbf L_{-A}\phi_1,\mathbf L_{-B}\phi_2,
                                      \mathbf L_{-C}\phi_3)
 \nonumber\\
 &\qquad=(-1)^{p_1A+p_2C}
 \left.\rho_a(\mathbf L_{-A}\phi_1,\mathbf L_{-B}\phi_2,
                                      \mathbf L_{-C}\phi_3)
                  \right|_{p_1=p_2=p_3=0}.
 \label{eq:ns-primary-phase}
\end{align}
The restriction on the right changes only the intrinsic primary parities.
It expresses the signs from moving odd supercurrent contours past those
primaries. The branch convention removes this dependence from the local
component values. Multiplying by the same sign undoes it. Consequently,
\begin{align}
 &\hat\rho_a\bigl(
 \Psi_{-\mathsf A}\mathbf1\otimes\mathbf L_{-A}\phi_1,
 \Psi_{-\mathsf B}\mathbf1\otimes\mathbf L_{-B}\phi_2,
 \Psi_{-\mathsf C}\mathbf1\otimes\mathbf L_{-C}\phi_3\bigr)
 \nonumber\\
 &\quad=(-1)^{B\mathsf C}(-1)^{p_1A+p_2C}
 \rho_{\F}(\Psi_{-\mathsf A}\mathbf1,
           \Psi_{-\mathsf B}\mathbf1,\Psi_{-\mathsf C}\mathbf1)
 \nonumber\\
 &\qquad\times
 \rho_a(\mathbf L_{-A}\phi_1,\mathbf L_{-B}\phi_2,
                                      \mathbf L_{-C}\phi_3).
 \label{eq:ns-vertex}
\end{align}
In particular, for even primaries only the middle-to-ket crossing remains.
This is the all-NS convention of the SCblock notes, with every argument
and both sources of its sign displayed.

At an NS--R--R vertex, let the physical arguments be
$(\mathbf L_{-A}\phi_1,\mathbb L_{-B}w_2^\alpha,
\mathbb L_{-C}w_3^\gamma)$. The contour normalization relative to the
physical form contributes
$(-i)^{A\bmod2}(-1)^{f(A+B+|\alpha|)}$.
The factor $-i$ accompanies odd NS supercurrent parity at infinity.
The $f$-dependent sign fixes the Ramond ground convention: when both
auxiliary grounds are $u^0$, the odd physical form requires
$\eta\mapsto(-1)^f\eta$. Starting from the ground values, an NS
supercurrent step changes this contour factor by $-i(-1)^{A+f}$,
and changing the second Ramond parity changes it by $(-1)^f$.
These are the phase ratios in the transported supercurrent Ward
identities. The auxiliary fermion contour at infinity has multiplier
$+i$, so the product with the physical $-i$ is one for the even field
$\psi G$ entering both Virasoro generators. This explains why the phase
must accompany the chosen Virasoro embedding.

For an odd NS primary the normalization conversion also contributes
$(-1)^{p_1\epsilon'_2}$ relative to the literal rearrangement
\eqref{eq:literal-moving}. Using
$\epsilon_1=A+p_1$ and
$\epsilon'_1+\epsilon'_2+\epsilon'_3=0$, the combined sign is
\begin{equation}
 \epsilon_1(\epsilon'_2+\epsilon'_3)+\epsilon_2\epsilon'_3
                 +p_1\epsilon'_2
 =A\epsilon'_1+(\epsilon_2+p_1)\epsilon'_3\pmod2.
 \label{eq:ramond-moving}
\end{equation}
The complete local definition is therefore
\begin{align}
 &\hat\rho_f^{(\eta)}\bigl(
 \Psi_{-\mathsf A}\mathbf1\otimes\mathbf L_{-A}\phi_1,
 \Psi_{-\mathsf B}u^{\mathsf b}\otimes\mathbb L_{-B}w_2^\alpha,
 \Psi_{-\mathsf C}u^{\mathsf c}\otimes\mathbb L_{-C}w_3^\gamma\bigr)
 \nonumber\\
 &\quad=(-i)^{A\bmod2}(-1)^{f(A+B+|\alpha|)}
        (-1)^{A\mathsf A+(B+|\alpha|+p_1)(\mathsf C+\mathsf c)}
 \nonumber\\
 &\qquad\times\rho_{\F}(\Psi_{-\mathsf A}\mathbf1,
                 \Psi_{-\mathsf B}u^{\mathsf b},
                 \Psi_{-\mathsf C}u^{\mathsf c})
 \rho_f^{(\eta)}(\mathbf L_{-A}\phi_1,
                     \mathbb L_{-B}w_2^\alpha,\mathbb L_{-C}w_3^\gamma).
 \label{eq:ramond-vertex}
\end{align}
Both sides vanish when the physical or auxiliary parity constraint fails.
The powers of $i$ always use a parity in $\{0,1\}$; the powers of $-1$
may be reduced modulo two.

\subsection{Sewing the enlarged block}

The enlarged PBW sum uses the physical inverse Gram entries in
\eqref{eq:physical-block}, the tensor states in
\eqref{eq:ns-vertex} and \eqref{eq:ramond-vertex}, and sums over both
physical and auxiliary internal labels. Its edge factor is obtained by
\begin{equation}
 q_e^{|A_{e,1}|}\eta_e^{\epsilon_e}
 \longmapsto q_e^{|A_{e,1}|+|\mathsf A_e|}
                         \eta_e^{\epsilon_e}(-\eta_e)^{\epsilon'_e},
 \qquad K(\epsilon)\longmapsto K(\epsilon+\epsilon').
 \label{eq:enlarged-block}
\end{equation}
The minus in $-\eta_e$ is exactly the inverse auxiliary Gram sign
\eqref{eq:aux-metric}; it is not a change of the physical spin structure.
Auxiliary words and ground labels agree at the two ends of an unsplit
edge. Replace each physical vertex by its fully specified enlarged form
above, and multiply by
\begin{equation}
 \prod_{v:\NS,\R,\R}i^{(A_{v,1}+\mathsf A_{v,1})\bmod2}.
 \label{eq:vertex-weight}
\end{equation}
This last factor defines the enlarged sewing convention used here. It
measures the total NS descendant parity, so it is constant on a double
Virasoro module. For theta the NS inverse Gram matrix identifies the two
physical parities and the auxiliary words agree, giving
$i^{(A+\mathsf A)\bmod2}i^{(A+\mathsf A)\bmod2}=(-1)^{A+\mathsf A}$,
as in the draft. On a general graph the vertex weight remains explicit;
it is a linear combination of parity sewings, rather than a single spin
lift. Equations \eqref{eq:physical-block}, \eqref{eq:enlarged-block},
and \eqref{eq:vertex-weight}, with the displayed local forms, specify
every factor and every summed index of $\widehat{\mathbf F}_{\hat\Gamma}$.

\subsection{Deriving the fermion factor}

To separate the two factors, combine the Ramond vertex with its sewing
weight first. The apparently complicated NS-slot signs cancel:
\begin{equation}
 i^{(A+\mathsf A)\bmod2}(-i)^{A\bmod2}(-1)^{A\mathsf A}
       =i^{\mathsf A\bmod2}.
 \label{eq:phase-cancellation}
\end{equation}
The remaining signs split into three parts:
\begin{align}
 (-1)^{f(A+B+|\alpha|)+(B+|\alpha|+p_1)(\mathsf C+\mathsf c)}
 &=\underbrace{(-1)^{f(1+\epsilon_3)}}_{\text{physical}}
   \underbrace{(-1)^{p_1\epsilon'_3}}_{\text{auxiliary}}
   \underbrace{(-1)^{\epsilon_2\epsilon'_3}}_{\text{mixed}},
 \label{eq:ramond-separation}\\
 (-1)^{B\mathsf C+p_1A+p_2C}
 &=\underbrace{(-1)^{p_1A+p_2C}}_{\text{physical}}
   \underbrace{(-1)^{p_2\epsilon'_3}}_{\text{auxiliary}}
   \underbrace{(-1)^{\epsilon_2\epsilon'_3}}_{\text{mixed}}.
 \label{eq:ns-separation}
\end{align}
The first line uses the nonzero-component condition
$A+B+C+|\alpha|+|\gamma|=f$ and
$\epsilon_3=C+|\gamma|$. Thus
$f(A+B+|\alpha|)=f(f+\epsilon_3)=f(1+\epsilon_3)$ modulo two.
The second line uses $\epsilon_2=B+p_2$. These identities remove the
need to track several separate Ramond signs in the graph contraction.

Collect the auxiliary factors just derived, together with their edge
weights and operator-reordering sign. This gives the following explicit
fermion block:
\begin{align}
 \mathbf F_{\F,\hat\Gamma}
 =\sum_{\{\mathsf A_e,\mathsf a_e\}}{}&(-1)^{K(\epsilon')}
     \prod_e q_e^{|\mathsf A_e|}(-\eta_e)^{\epsilon'_e}
 \nonumber\\
 &\times\prod_{v:\NS,\NS,\NS}
 \Bigl[(-1)^{p_{v,2}\mathsf A_{v,3}}
 \rho_{\F}(\Psi_{-\mathsf A_{v,1}}\mathbf1,
           \Psi_{-\mathsf A_{v,2}}\mathbf1,
           \Psi_{-\mathsf A_{v,3}}\mathbf1)\Bigr]
 \nonumber\\
 &\times\prod_{v:\NS,\R,\R}
 \Bigl[i^{\mathsf A_{v,1}\bmod2}
       (-1)^{p_{v,1}(\mathsf A_{v,3}+\mathsf a_{v,3})}
 \rho_{\F}(\Psi_{-\mathsf A_{v,1}}\mathbf1,
           \Psi_{-\mathsf A_{v,2}}u^{\mathsf a_{v,2}},
           \Psi_{-\mathsf A_{v,3}}u^{\mathsf a_{v,3}})\Bigr].
 \label{eq:fermion-block}
\end{align}
There is one independent auxiliary word on each internal edge and one
additional binary ground label on each Ramond edge; these are the summed
indices. The labels $\mathsf A_{v,i},\mathsf a_{v,i}$ are the labels of
the edge incident at that slot. External labels are held fixed. Terms
with $\sum_i\epsilon'_{v,i}\ne0$ vanish. The signs $p_{v,i}$ are fixed
physical primary parities, not extra fermion summation variables.
For recovery use $\mathbf1$ at each original external NS puncture and
$u^0$ at each original external R puncture. This choice normalizes open
Ramond paths to one at level zero and places no restriction on the physical
external state.

In the NS--R--R theta channel this formula provides an immediate check on
all three sources of signs. Auxiliary parity conservation gives
$\epsilon'_1+\epsilon'_2+\epsilon'_3=0$, so the three edge-metric signs
multiply to one. The two local primary-parity signs cancel, while the
two factors $i^{\mathsf A\bmod2}$ give $(-1)^{\mathsf A}$. Hence
\begin{align}
 \mathbf F_{\F,\theta}
 =\sum_{\substack{\mathsf A,\mathsf B,\mathsf C,\mathsf b,\mathsf c\\
       \mathsf A+\mathsf B+\mathsf C+\mathsf b+\mathsf c=0\ (\mathrm{mod}\ 2)}}{}
 &(-1)^{\mathsf A+\epsilon'_1\epsilon'_2+
                   \epsilon'_1\epsilon'_3+\epsilon'_2\epsilon'_3}
 q_1^{|\mathsf A|}q_2^{|\mathsf B|}q_3^{|\mathsf C|}
 \eta_1^{\epsilon'_1}\eta_2^{\epsilon'_2}\eta_3^{\epsilon'_3}
 \nonumber\\
 &\times\bigl[\rho_{\F}(\Psi_{-\mathsf A}\mathbf1,
           \Psi_{-\mathsf B}u^{\mathsf b},
           \Psi_{-\mathsf C}u^{\mathsf c})\bigr]^2,
 \label{eq:fermion-theta}
\end{align}
where $(\epsilon'_1,\epsilon'_2,\epsilon'_3)
=(\mathsf A,\mathsf B+\mathsf b,\mathsf C+\mathsf c)$ modulo two.
This is the fermion factor in the draft.
\section{Convolution and recovery of the physical parity convention}

A term in the physical block has powers $r_e$ and parities $\epsilon_e$;
a term in the auxiliary block has powers $s_e$ and parities $\epsilon'_e$.
Their product adds levels on each edge and adds parities modulo two.
There are only two sources of a mixed sign. At a reversed pair of
punctures, subtracting the separate physical and auxiliary signs from the
tensor-state sign leaves
\begin{equation}
 (\epsilon_{v,i}+\epsilon'_{v,i})(\epsilon_{w,j}+\epsilon'_{w,j})
 -\epsilon_{v,i}\epsilon_{w,j}-\epsilon'_{v,i}\epsilon'_{w,j}
 =\epsilon_{v,i}\epsilon'_{w,j}+\epsilon'_{v,i}\epsilon_{w,j}.
 \label{eq:mixed-crossing}
\end{equation}
Summing over the reversed pairs gives
$K(\epsilon+\epsilon')-K(\epsilon)-K(\epsilon')$.
The local vertex identities \eqref{eq:ramond-separation} and
\eqref{eq:ns-separation} supply the remaining
$\sum_v\epsilon_{v,2}\epsilon'_{v,3}$. Thus the required convolution is
\begin{align}
 \left(\prod_e q_e^{r_e}\eta_e^{\epsilon_e}\right)\star
 \left(\prod_e q_e^{s_e}\eta_e^{\epsilon'_e}\right)
 ={}&(-1)^{K(\epsilon+\epsilon')-K(\epsilon)-K(\epsilon')
                    +\sum_v\epsilon_{v,2}\epsilon'_{v,3}}
 \prod_e q_e^{r_e+s_e}\eta_e^{\epsilon_e+\epsilon'_e}.
 \label{eq:star}
\end{align}
Internal edge parities determine the corresponding vertex-slot
parities. External parities are fixed labels of the two inputs; the
result carries their sum. The sign is
bilinear in the two parity assignments. This proves associativity by
expanding either parenthesization of a product of three monomials.
Keep the physical factor on the left, since the product need not be
commutative. The $\eta_e$ are formal parity labels with ordinary
relation $\eta_e^2=1$; a convolution square may carry an additional
minus sign from \eqref{eq:star}. Evaluate the final physical block at
$\eta_e=\pm1$ after the coefficient calculation. Substitution before
convolution would lose these signs.

The purely physical factors in \eqref{eq:ramond-separation} and
\eqref{eq:ns-separation} remain to be included. Their product, written
entirely in terms of the retained parities, is
\begin{equation}
 (-1)^{\sum_{v:\NS,\R,\R}f_v(1+\epsilon_{v,3})
       +\sum_{v:\NS,\NS,\NS}
       [p_{v,1}(\epsilon_{v,1}+p_{v,1})+
        p_{v,2}(\epsilon_{v,3}+p_{v,3})]}.
 \label{eq:transport}
\end{equation}
No descendant word is needed to evaluate it. For even NS primaries only
$(-1)^{\sum_{v:\NS,\R,\R}f_v(1+\epsilon_{v,3})}$ remains.
This is a fixed conversion between the physical and enlarged conventions,
not another sum over spin structures.

For any series, $[\prod_e q_e^{r_e}\eta_e^{\epsilon_e}]$ denotes its
scalar coefficient, whereas $[\prod_e q_e^{r_e}]$ retains the entire
parity polynomial at those levels. With the vertex labels and external
states held fixed, the full relation is
\begin{align}
 \widehat{\mathbf F}_{\hat\Gamma}
 =\sum_{\{r,\epsilon\}}{}&
 (-1)^{\sum_{v:\NS,\R,\R}f_v(1+\epsilon_{v,3})
       +\sum_{v:\NS,\NS,\NS}
        [p_{v,1}(\epsilon_{v,1}+p_{v,1})+
         p_{v,2}(\epsilon_{v,3}+p_{v,3})]}
 \nonumber\\
 &\times\mathbf F_{\hat\Gamma}
       \left[\prod_e q_e^{r_e}\eta_e^{\epsilon_e}\right]
       \left(\prod_e q_e^{r_e}\eta_e^{\epsilon_e}\right)
                                  \star\mathbf F_{\F,\hat\Gamma}.
 \label{eq:convolution}
\end{align}
The sum is over the physical levels and parity components allowed by
\eqref{eq:vertex-selection}; equivalently one may include all components
and set forbidden coefficients to zero. The displayed sign is applied
once per physical coefficient. After solving the convolution it is undone
by multiplying by the same sign. In theta the equal vertex labels make
these signs cancel in pairs, so
$\widehat{\mathbf F}_{\theta}=\mathbf F_{\theta}\star\mathbf F_{\F,\theta}$.
On other graphs the conversion must be retained. The preceding derivation
identifies every factor with a PBW edge or vertex contribution, rather
than assuming a factorization of the final numerical series.
\section{Reduction to double Virasoro}

Each enlarged internal module is a direct sum of modules for the two
commuting Virasoro algebras. Their primaries are the $v_n$ and
$v_n^\alpha$ of the draft. The label ranges and parameters are
\begin{gather}
 n\in\tfrac12\mathbb Z\quad(\NS),\qquad
 n\in\tfrac12\mathbb Z+\tfrac14\quad(\R),\qquad
 \alpha=0,1,\nonumber\\
 Q=b+b^{-1},\qquad c=\tfrac32+3Q^2,\qquad b\ne0,\pm1,
 \nonumber\\
 c^{(1)}=1+\frac{3Q^2}{1-b^2},\qquad
 c^{(2)}=1+\frac{3Q^2}{1-b^{-2}},\nonumber\\
 h_{e,n}^{(1)}=\frac{Q^2/4-(P_e+2nb)^2}{2(1-b^2)},\qquad
 h_{e,n}^{(2)}=\frac{Q^2/4-(P_e+2n/b)^2}{2(1-b^{-2})}.
 \label{eq:weights}
\end{gather}
For an NS module $P_e^2=Q^2/4-2h_e$; for a Ramond module
$P_e^2=2\beta_e^2$ and $h_e=c/24-\beta_e^2$.
The label on $v_n^0,v_n^1$ is the total even or odd parity of an enlarged
primary, which mixes auxiliary and physical Ramond grounds. The states
$w^\pm$ used in the PBW sum are the physical grounds before this change
of basis.

Subtracting the enlarged ground weight gives the starting level of each
branch. The NS auxiliary ground has weight zero, and the Ramond auxiliary
ground has weight $1/16$. Therefore
\begin{equation}
 \begin{array}{c|c|c}
 &\text{branch starting level}&\text{total parity }\epsilon_e\\ \hline
 \NS& h_{e,n}^{(1)}+h_{e,n}^{(2)}-h_e=2n^2
                                      &(2n+p_e)\bmod2\\
 \R& h_{e,n}^{(1)}+h_{e,n}^{(2)}-h_e-\tfrac1{16}=2n^2-\tfrac18
                                      &\alpha_e
 \end{array}
 \label{eq:branches}
\end{equation}
Thus $\epsilon$ in this section is the parity of the sewn branch state;
in the PBW convolution it was the parity of the physical factor.
All states within a fixed double-Virasoro module have the same parity.

The two types of normalized branching coefficient are
\begin{align}
 \mathbf B_a(n_1,n_2,n_3)
 &=\frac{\hat\rho_a(v_{1,n_1},v_{2,n_2},v_{3,n_3})}
         {\|v_{1,n_1}\|\|v_{2,n_2}\|\|v_{3,n_3}\|},\nonumber\\
 \mathbb B_f^{(\eta)}(n_1,n_2,n_3;\alpha_2,\alpha_3)
 &=\frac{\hat\rho_f^{(\eta)}
             (v_{1,n_1},v_{2,n_2}^{\alpha_2},v_{3,n_3}^{\alpha_3})}
         {\|v_{1,n_1}\|\|v_{2,n_2}^{\alpha_2}\|\|v_{3,n_3}^{\alpha_3}\|}.
 \label{eq:branching-definitions}
\end{align}
Here $\|v\|^2=\langle\!\langle v|v\rangle$ is a bilinear BPZ norm,
not an absolute square. Choose its roots consistently for all vertices
and edges. We suppress only the fixed representation parameters in
$\mathbf B,\mathbb B$; the branch labels and Ramond parities are displayed.

First take normalized branch primaries at the external punctures. The
enlarged block is then
\begin{align}
 \widehat{\mathbf F}_{\hat\Gamma}
 =\sum_{\{n_e,\alpha_e\}}{}&(-1)^{K(\epsilon)}
 \prod_{e:\NS}q_e^{2n_e^2}\eta_e^{2n_e+p_e}
 \prod_{e:\R}q_e^{2n_e^2-1/8}\eta_e^{\alpha_e}
 \prod_{v:\NS,\R,\R}i^{(2n_{v,1})\bmod2}
 \nonumber\\
 &\times\prod_{v:\NS,\NS,\NS}
             \mathbf B_{a_v}(n_{v,1},n_{v,2},n_{v,3})
 \nonumber\\
 &\times\prod_{v:\NS,\R,\R}
             \mathbb B_{f_v}^{(\eta_v)}
                  (n_{v,1},n_{v,2},n_{v,3};\alpha_{v,2},\alpha_{v,3})
 \prod_{j=1}^2F_{\hat\Gamma}(\{h_{e,n_e}^{(j)}\};c^{(j)}).
 \label{eq:double-virasoro}
\end{align}
The sum and edge products run over internal edges; external branch labels
are fixed. The argument list of each Virasoro block includes both the
internal weights and the fixed external weights. A slot label $n_{v,i}$
is the label of its incident edge, so a self-loop has one branching sum
shared by its two slots. Keep only the local parity components allowed by
\eqref{eq:vertex-selection}, now applied to total branch parities.

To see the reduction, change basis separately on every internal edge.
The BPZ form becomes the product of the two Virasoro Gram matrices and
the primary norm. The two adjacent normalized branching coefficients
supply the inverse primary norm once per edge. The Virasoro Ward
identities then reduce descendants independently in the two copies.
For example, at an NS--R--R vertex,
\begin{align}
 &\hat\rho_f^{(\eta)}
 (L_{-1}^{(j)}v_{1,n_1},v_{2,n_2}^{\alpha_2},v_{3,n_3}^{\alpha_3})
 \nonumber\\
 &\qquad=(h_{1,n_1}^{(j)}+h_{2,n_2}^{(j)}-h_{3,n_3}^{(j)})
 \hat\rho_f^{(\eta)}
 (v_{1,n_1},v_{2,n_2}^{\alpha_2},v_{3,n_3}^{\alpha_3}),
 \qquad j=1,2.
 \label{eq:local-virasoro-ward}
\end{align}
This identity uses the full local form \eqref{eq:ramond-vertex}; deleting
its contour phase while keeping the embeddings would spoil the reduction.
The sewing weight \eqref{eq:vertex-weight} is constant on each branch
and becomes $i^{(2n_{v,1})\bmod2}$ in \eqref{eq:double-virasoro}.
An external tensor-PBW state is handled by the finite expansion in
Sec.~\ref{sec:external}, followed by the same Ward reduction.

In computation, evaluate the Virasoro factors by CCY $c$-recursion.
Its large-$c$ seed contains the $L_{-1}$ global block and the universal
Schottky vacuum factor. Store pole geometry, fusion polynomials, residue
normalizations, global vertices, and repeated recursion states. At
NS--R--R vertices obtain the outer branching coefficients recursively
from $n_{\NS}=0,\pm\tfrac12$ and
$n_{\R}=\pm\tfrac14,\pm\tfrac34$, reusing intermediate coefficients
and the coefficients of the $L_{\pm1}$ descendant expansions.
When a Ward equation has a zero leading coefficient, use the alternative
with $L_1$ at a Ramond slot. The validation driver obtains its all-NS
coefficients from free-field primary vectors and caches those local
values. The remaining Virasoro cutoff depends on the branch: subtract
its starting level in \eqref{eq:branches} before evaluating descendants.
This prevents computing blocks beyond the requested physical degree.

\section{Ramond loops, cancellation, and the constant term}

Every vertex meets zero or two Ramond edges. The Ramond subgraph therefore
consists of disjoint closed cycles and paths with external Ramond ends.
A self-loop is a closed cycle with one vertex. Let $C$ be any closed
Ramond cycle. We write $K(C)$ for the already defined inversion count
$K$, evaluated with parity one on that cycle and zero on all other
internal and external punctures. No new sign rule is involved.

The physical coefficients on the cycle are related by
\begin{equation}
 \left(\prod_{e\in C}\eta_e\right)\star\mathbf F_{\hat\Gamma}
 =(-1)^{K(C)}\left(\prod_{v\in C}i\eta_v\right)
                                      \mathbf F_{\hat\Gamma}.
 \label{eq:physical-sector}
\end{equation}
The left side flips the retained parity on every edge of $C$, with the
convolution signs of \eqref{eq:star}. The right side fixes the relative
coefficients of those components. It is a sector identity for the formal
series; it does not assign complex values to the spin lifts.

Locally this identity follows from the odd ground map
$w^+\mapsto e^{-i\pi/4}w^-$ and
$w^-\mapsto-e^{i\pi/4}w^+$, extended to a Ramond PBW word with the
factor $(-1)^A$. This map preserves the ground BPZ form
$\operatorname{diag}(1,i)$. Acting on both Ramond slots multiplies the
physical form by $-i\eta_v(-1)^{\epsilon_{v,2}}$; the ground values
\eqref{eq:physical-vertices} verify it, and the supercurrent Ward
identities extend it to descendants. Combining these local identities
around the cycle and reordering the odd maps gives
\eqref{eq:physical-sector}. Crossing the NS slot twice contributes
$(-1)^{2\epsilon_{v,1}}=1$, including for an odd NS primary.

The cycle monomial commutes with admissible factors of the convolution:
its commutator exponent is
\begin{equation}
 \sum_{v\in C}(\epsilon_{v,2}+\epsilon_{v,3})
                  =2\sum_{e\in C}\epsilon_e=0\pmod2.
 \label{eq:loop-centrality}
\end{equation}
Each Ramond edge is counted once at each end. The global terms from
$K$ cancel in the commutator because they are symmetric in its two inputs.
Multiplying physical coefficients by \eqref{eq:transport} changes the
cycle eigenvalue in \eqref{eq:physical-sector} by
$(-1)^{\sum_{v\in C}f_v}$: the parity in slot 3 is flipped once at each
vertex of $C$. Applying the same ground-map argument to the auxiliary
form gives
\begin{equation}
 \left(\prod_{e\in C}\eta_e\right)\star\mathbf F_{\F,\hat\Gamma}
 =(-1)^{K(C)+\sum_{v\in C}p_{v,1}}
       \left(\prod_{v\in C}i\right)\mathbf F_{\F,\hat\Gamma}.
 \label{eq:aux-sector}
\end{equation}
Move this commuting monomial between the two factors of
\eqref{eq:convolution} and cancel the common phase. In one line,
\begin{equation}
 \left[(-1)^{\sum_{v\in C}f_v}\prod_{v\in C}\eta_v
                  -(-1)^{\sum_{v\in C}p_{v,1}}\right]
                               \widehat{\mathbf F}_{\hat\Gamma}=0.
 \label{eq:loop-cancellation}
\end{equation}
Thus the enlarged block vanishes whenever
\begin{equation}
 (-1)^{\sum_{v\in C}(f_v+p_{v,1})}\prod_{v\in C}\eta_v=-1.
 \label{eq:odd-cycle}
\end{equation}
Its two factors belong to incompatible cycle sectors; the physical block
itself can be nonzero. In theta the two equal $f$ labels and primary
parities cancel, leaving $\eta\eta'=-1$.

The auxiliary constant term makes both the obstruction and the restricted
inverse concrete. With the external ground states chosen in
Sec.~\ref{sec:local-signs}, all NS auxiliary slots are vacuum at level zero. Around a closed
Ramond cycle the auxiliary grounds must then be either all $u^0$ or all
$u^1$. The first choice contributes one. The second contributes $i$ from
each vertex, $-1$ from each inverse ground metric, its spin lifts, and the
ordering and primary-parity signs. Consequently,
\begin{equation}
 \left.\mathbf F_{\F,\hat\Gamma}\right|_{q_e=0}
 =\prod_{C}^{\star}
 \left[1+(-1)^{K(C)+\sum_{v\in C}p_{v,1}}
              \left(\prod_{v\in C}(-i)\right)
              \left(\prod_{e\in C}\eta_e\right)\right].
 \label{eq:constant-projectors}
\end{equation}
The product runs over the closed Ramond cycles and is taken using $\star$;
its factors commute. The empty product is one, and open paths with
$u^0$ ends contribute one. Each term after the $1$ in a bracket has
convolution square one: the cycle monomial has square
$(-1)^{\text{number of vertices on }C}$, canceled by the squared
factors $-i$. Hence each bracket acts by two on its matching sector and
by zero on the other sector. The constant is invertible on the matching
simultaneous sector, but not on the full coefficient space if a closed
Ramond cycle is present.

\section{Splitting an obstructed edge and recovering the block}

\subsection{The inserted sphere and its PBW matrix element}

Choose one edge on each cycle satisfying \eqref{eq:odd-cycle}. On a chosen
edge $e$, insert a new sphere with coordinate $w$, sew its punctures
$0,\infty$ to the original ends, and place $\Theta\psi$ at $w=1$.
If $z_L,z_R$ are the original local coordinates, the two sewing equations
are $z_Lw=q_{e_L}$ and $z_R/w=q_{e_R}$. Multiplying them gives
\begin{equation}
 q_e=q_{e_L}q_{e_R},\qquad \eta_e=\eta_{e_L}\eta_{e_R}.
 \label{eq:split-map}
\end{equation}
The spin lifts multiply in the aligned Ramond frames, with no extra odd
insertion sign because $\Theta\psi$ is even. The operator is defined by
\begin{equation}
 \Theta u^{\mathsf a}=(-1)^{\mathsf a+1}u^{1-\mathsf a},
 \qquad\{\Theta,\psi_r\}=\{\Theta,G_r\}=0,
 \qquad[\Theta,L_n]=0.
 \label{eq:theta}
\end{equation}
Both factors anticommute with physical $G_r$, so their product commutes
with the physical SCA. Also $\Theta$ commutes with $L^{(1,2)}$:
it commutes with the even fermion bilinears and with $\psi G$.
The inserted vertex is a radial matrix element, without the additional
outer-vertex weight \eqref{eq:vertex-weight}.

The physical descendant words at the inserted sphere are $B_L,B_R$.
Their auxiliary words are $\mathsf B_L,\mathsf B_R$, with ground labels
$\mathsf b_L,\mathsf b_R$. The matrix element is
\begin{align}
 &\Big\langle\!\Big\langle
 \Psi_{-\mathsf B_R}u^{\mathsf b_R}\otimes\mathbb L_{-B_R}w^{\alpha_R}
 \Big|\Theta\psi(1)\Big|
 \Psi_{-\mathsf B_L}u^{\mathsf b_L}\otimes\mathbb L_{-B_L}w^{\alpha_L}
 \Big\rangle
 \nonumber\\
 &\qquad=\langle\!\langle\Psi_{-\mathsf B_R}u^{\mathsf b_R}|
                 \Theta\psi(1)|\Psi_{-\mathsf B_L}u^{\mathsf b_L}\rangle
                (\mathbb G_{\beta_e})_{B_R\alpha_R,B_L\alpha_L}.
 \label{eq:split-pbw}
\end{align}
The order fixes $\Theta$ next to the outgoing bra. Each new segment has
its own physical inverse Gram matrix, auxiliary inverse metric, and
factor $q^{|B|+|\mathsf B|}\eta^{\epsilon}(-\eta)^{\epsilon'}$ as in
\eqref{eq:enlarged-block}. All endpoint indices are summed as before.
Using the symmetry of the physical BPZ Gram matrix, its contraction reduces to
\begin{align}
 &\sum_{B_L,B_R,\alpha'_L,\alpha'_R}
 \mathbb G_{\beta_e}^{A_L\alpha_L,B_L\alpha'_L}
 (\mathbb G_{\beta_e})_{B_L\alpha'_L,B_R\alpha'_R}
 \mathbb G_{\beta_e}^{B_R\alpha'_R,A_R\alpha_R}
 =\mathbb G_{\beta_e}^{A_L\alpha_L,A_R\alpha_R}.
 \label{eq:split-gram}
\end{align}
This contracts the identity operator on the physical line. Its levels and
parities agree on the two segments. Auxiliary parities agree because
$\Theta\psi$ is even, but auxiliary levels may differ through the
nonzero modes of $\psi(1)$. Using only the auxiliary matrix element in
\eqref{eq:split-pbw} defines
$\mathbf F_{\F,\hat\Gamma}[\Theta\psi]$ with the same outer vertices
as \eqref{eq:fermion-block}.

\subsection{The inserted double-Virasoro sum}

The field $\psi$ is proportional to the NS branch primary
$v_{1/2}$ at $P=Q/2$. Its two Virasoro weights are
\begin{equation}
 h_\psi^{(1)}=-\frac{1+2b^2}{2(1-b^2)},\qquad
 h_\psi^{(2)}=\frac{b^2+2}{2(1-b^2)},\qquad
 h_\psi^{(1)}+h_\psi^{(2)}=\tfrac12.
 \label{eq:psi-weights}
\end{equation}
Their degenerate fusion rules give $n_{e_R}=n_{e_L}\pm\tfrac12$.
The odd $\psi$ reverses the total Ramond parity and $\Theta$ reverses it
back, so the two segments carry the same $\alpha_e$. The normalized
middle coefficient is
\begin{equation}
 \mathbb B(n_{e_R},n_{e_L},\alpha_e)
 =\frac{\langle\!\langle v_{e,n_{e_R}}^{\alpha_e}|
                    \Theta\psi(1)|v_{e,n_{e_L}}^{\alpha_e}\rangle}
        {\|v_{e,n_{e_R}}^{\alpha_e}\|\|v_{e,n_{e_L}}^{\alpha_e}\|}.
 \label{eq:middle-branching}
\end{equation}
It is zero outside the stated fusion rule. Since $\Theta$ commutes with
both Virasoro copies, the descendant dependence is entirely in their
ordinary punctured blocks.

With normalized branch primaries at the original external punctures,
the full sum is
\begin{align}
 \widehat{\mathbf F}_{\hat\Gamma}[\Theta\psi]
 =\sum_{\{n,\alpha\}}{}&(-1)^{K(\epsilon)}
 \prod_{v:\NS,\R,\R}i^{(2n_{v,1})\bmod2}
 \prod_{e:\NS}q_e^{2n_e^2}\eta_e^{2n_e+p_e}
 \prod_{\substack{e:\R\\e\text{ unsplit}}}
                       q_e^{2n_e^2-1/8}\eta_e^{\alpha_e}
 \nonumber\\
 &\times\prod_{e\text{ split}}
       q_{e_L}^{2n_{e_L}^2-1/8}q_{e_R}^{2n_{e_R}^2-1/8}
       (\eta_{e_L}\eta_{e_R})^{\alpha_e}
       \mathbb B(n_{e_R},n_{e_L},\alpha_e)
 \nonumber\\
 &\times\prod_{v:\NS,\NS,\NS}
               \mathbf B_{a_v}(n_{v,1},n_{v,2},n_{v,3})
 \nonumber\\
 &\times\prod_{v:\NS,\R,\R}
               \mathbb B_{f_v}^{(\eta_v)}
                  (n_{v,1},n_{v,2},n_{v,3};\alpha_{v,2},\alpha_{v,3})
 \nonumber\\
 &\times\prod_{j=1}^2F_{\hat\Gamma}[\psi]
                   (\{h_{e,n_e}^{(j)}\};\{h_\psi^{(j)}\};c^{(j)}).
 \label{eq:split-double-virasoro}
\end{align}
Vertex products here include only the original spheres; the middle
coefficients account for the new spheres once each. The internal branch
sum includes both labels $n_{e_L},n_{e_R}$ on a split edge, constrained
by the fusion rule, and one common parity $\alpha_e$. In an outer
branching coefficient, $n_{v,i}$ is the label of the segment adjacent to
that original vertex. The function $F_{\hat\Gamma}[\psi]$ denotes the
ordinary Virasoro block on this expanded channel, with those segment
weights, the original external weights, and one additional external weight
$h_\psi^{(j)}$ for each inserted sphere. It is computed by the same CCY
recursion. The original graph's $K$ applies because the inserted operator
is even and the segment parities agree. Outer and middle coefficients
use the same BPZ norm roots.

\subsection{Convolution before and after the diagonal restriction}

Equation \eqref{eq:split-gram} shows that the full physical monomial is
replaced by
\begin{equation}
 \prod_e q_e^{r_e}\eta_e^{\epsilon_e}
 \longmapsto
 \prod_{e\text{ unsplit}}q_e^{r_e}\eta_e^{\epsilon_e}
 \prod_{e\text{ split}}
       (q_{e_L}q_{e_R})^{r_e}(\eta_{e_L}\eta_{e_R})^{\epsilon_e}.
 \label{eq:split-monomial}
\end{equation}
Thus the complete unprojected convolution is
\eqref{eq:convolution} with
$\widehat{\mathbf F}_{\hat\Gamma}$ and
$\mathbf F_{\F,\hat\Gamma}$ replaced by their inserted blocks, and
every physical monomial replaced by \eqref{eq:split-monomial}. The
physical coefficient itself is still that of the original graph.
This is an exact substitution in the displayed formula, including its
physical correction and mixed sign. In particular, a physical power
$r_e$ and auxiliary powers $(s_{e_L},s_{e_R})$ produce
$(r_e+s_{e_L},r_e+s_{e_R})$; unequal split powers remain included.

To recover the physical block it suffices to keep the terms with equal
powers on each split pair. Define the subscript $\mathrm{diag}$ precisely
by coefficient extraction:
\begin{align}
 &\widehat{\mathbf F}_{\hat\Gamma}[\Theta\psi]_{\mathrm{diag}}
                  \left[\prod_e q_e^{r_e}\eta_e^{\epsilon_e}\right]
 \nonumber\\
 &\quad=\widehat{\mathbf F}_{\hat\Gamma}[\Theta\psi]
 \left[\prod_{e\text{ unsplit}}q_e^{r_e}\eta_e^{\epsilon_e}
       \prod_{e\text{ split}}
          (q_{e_L}q_{e_R})^{r_e}(\eta_{e_L}\eta_{e_R})^{\epsilon_e}
 \right].
 \label{eq:diagonal-definition}
\end{align}
Use the identical definition for the fermion factor. Since the physical
powers on a split pair already agree, equality of the final powers is
equivalent to equality of the auxiliary powers. Hence the diagonal
restriction commutes with the convolution.

On the fermion side this keeps $\Theta\psi_0$. Its action is particularly
simple:
\begin{equation}
 \Theta\psi_0\Psi_{-\mathsf A}u^{\mathsf a}
   =\Psi_{-\mathsf A}\Theta\psi_0u^{\mathsf a}
   =\frac{(-1)^{\mathsf a}}{\sqrt2}\Psi_{-\mathsf A}u^{\mathsf a}.
 \label{eq:zero-mode}
\end{equation}
Moving each of $\Theta$ and $\psi_0$ through the nonzero fermions gives
$(-1)^{\mathsf A}$, so those two signs cancel. Only the ground label
remains. The ground factor on the affected cycle changes from the bracket
in \eqref{eq:constant-projectors} to
\begin{equation}
 \frac1{\sqrt2}
 \left[1-(-1)^{K(C)+\sum_{v\in C}p_{v,1}}
              \left(\prod_{v\in C}(-i)\right)
              \left(\prod_{e\in C}\eta_e\right)\right].
 \label{eq:inserted-ground}
\end{equation}
It now selects the previously missing sector. Each closed cycle acts by
two on its selected sector, with an additional $1/\sqrt2$ for each
insertion. No inverse on the unrestricted coefficient space is asserted.

The Virasoro calculation can enforce the same restriction before summing:
\begin{equation}
 2n_{e_L}^2-\tfrac18+k_{e_L}^{(1)}+k_{e_L}^{(2)}
 =2n_{e_R}^2-\tfrac18+k_{e_R}^{(1)}+k_{e_R}^{(2)}.
 \label{eq:diagonal-cutoff}
\end{equation}
Here $k^{(1)},k^{(2)}$ are the nonnegative integer descendant levels in
the two copies. The restriction concerns their combined power together
with the branch starting level. It does not set $n_{e_L}=n_{e_R}$ or
match the descendant levels separately in each copy. Use it together with
the physical cutoff to discard unneeded branch tuples and recursion
states, and reuse the middle and outer branching coefficients.

\subsection{The coefficient recursion}

After inserting once on every obstructed cycle, the diagonal auxiliary
constant is nonzero on the required sector. The following coefficient
equation makes the triangular solve explicit without introducing a new
name for a rephased physical block. The level indices $r_e,s_e,t_e$ are
nonnegative half-integers on NS edges and integers on Ramond edges;
coefficients absent from the corresponding series are zero. Here
$\prod_C2$ runs over the original closed Ramond cycles, and the other
constant factor runs over the edges split by insertions:
\begin{align}
 &\widehat{\mathbf F}_{\hat\Gamma}[\Theta\psi]_{\mathrm{diag}}
                    \left[\prod_e q_e^{r_e}\right]
 \nonumber\\
 &\quad=\sum_{\epsilon}
 (-1)^{\sum_{v:\NS,\R,\R}f_v(1+\epsilon_{v,3})
       +\sum_{v:\NS,\NS,\NS}
        [p_{v,1}(\epsilon_{v,1}+p_{v,1})+
         p_{v,2}(\epsilon_{v,3}+p_{v,3})]}
 \nonumber\\
 &\qquad\times\Bigg\{
     \left(\prod_C2\right)\left(\prod_{e\text{ split}}\frac1{\sqrt2}\right)
     \mathbf F_{\hat\Gamma}
                    \left[\prod_e q_e^{r_e}\eta_e^{\epsilon_e}\right]
     \prod_e\eta_e^{\epsilon_e}
 \nonumber\\
 &\qquad\quad+
 \sum_{\substack{s_e+t_e=r_e\ \text{for all }e\\s_e,t_e\ge0,\ \sum_e t_e>0}}
     \mathbf F_{\hat\Gamma}
                    \left[\prod_e q_e^{s_e}\eta_e^{\epsilon_e}\right]
     \left(\prod_e\eta_e^{\epsilon_e}\right)\star
     \mathbf F_{\F,\hat\Gamma}[\Theta\psi]_{\mathrm{diag}}
                    \left[\prod_e q_e^{t_e}\right]\Bigg\}.
 \label{eq:recovery}
\end{align}
This is an equality of parity polynomials. The second sum uses only lower
total physical levels because $\sum_e t_e>0$. Subtract it from the known
enlarged coefficient, divide by the displayed constant, and multiply each
parity component by \eqref{eq:transport} to obtain the desired physical
coefficient. At level zero the second sum is empty. For an unobstructed
calculation omit the insertion brackets, the diagonal restriction, and
the product over split edges.

The equations above use the canonical fermion normalization
\eqref{eq:aux-metric}. The implementation inserts $Q\Theta\psi$ in both
the enlarged numerator and the auxiliary factor. For that convention,
replace each $1/\sqrt2$ in the constant by $Q/\sqrt2$ and multiply each
middle coefficient by $Q$. The factors cancel from the recovered block.
At no stage is the physical PBW result an input to this recursion; it
enters only as an independent reference in the validation section.

\section{Theta, glasses, and Mercedes}

\subsection{Theta}

Use the two orders given after \eqref{eq:graph-sign}. The resulting signs
are
\begin{equation}
 K(\epsilon)=\epsilon_1\epsilon_2+\epsilon_1\epsilon_3+
              \epsilon_2\epsilon_3,
 \qquad
 (-1)^{\sum_{i<j}(\epsilon_i\epsilon'_j+\epsilon'_i\epsilon_j)}
 \quad\text{in the convolution}.
 \label{eq:theta-specialization}
\end{equation}
The two local mixed terms in \eqref{eq:star} cancel. The physical
correction \eqref{eq:transport} is one because the two vertices have equal
labels and share their corresponding parities. The cycle identity and
auxiliary ground factors become
\begin{align}
 (\eta_2\eta_3)\star\mathbf F_\theta
    &=\eta\eta'\,\mathbf F_\theta,\nonumber\\
 \left.\mathbf F_{\F,\theta}\right|_{q=0}
    &=1+\eta_2\eta_3,\qquad
 \left.\mathbf F_{\F,\theta}[\Theta\psi]_{\mathrm{diag}}\right|_{q=0}
    =\frac{1-\eta_2\eta_3}{\sqrt2}.
 \label{eq:theta-ground}
\end{align}
The ordinary factor selects $\eta\eta'=1$; the inserted factor selects
$\eta\eta'=-1$. Splitting either Ramond edge supplies the missing sector.
The local and sewing definitions reduce to the theta formulas in the draft.

\subsection{Glasses}

Use edge 1 as the NS bridge and edges 2,3 as Ramond self-loops, with vertex
triples $(1_L,2_L,2_R)$ and $(1_R,3_L,3_R)$. Take an even bridge primary.
Parity conservation fixes the same $f$ at both vertices. Now
\begin{equation}
 K(\epsilon)=0,
 \qquad (-1)^{\epsilon_2\epsilon'_2+\epsilon_3\epsilon'_3}
                 \quad\text{in the convolution},
 \qquad (-1)^{f(\epsilon_2+\epsilon_3)}
                 \quad\text{for \eqref{eq:transport}}.
 \label{eq:glasses-specialization}
\end{equation}
In particular $\eta_2\star\eta_2=\eta_3\star\eta_3=-1$.
The physical cycle identities are
\begin{equation}
 \eta_2\star\mathbf F_{\mathrm{gl}}
       =i\eta_{v_1}\mathbf F_{\mathrm{gl}},\qquad
 \eta_3\star\mathbf F_{\mathrm{gl}}
       =i\eta_{v_2}\mathbf F_{\mathrm{gl}}.
 \label{eq:glasses-sector}
\end{equation}
The ordinary auxiliary constant is $(1-i\eta_2)(1-i\eta_3)$.
Inserting on edge 2 replaces its factor by $(1+i\eta_2)/\sqrt2$, and
similarly for edge 3. The loop at vertex $v$ needs an insertion when
$(-1)^f\eta_v=-1$. Thus for $f=1$ the obstructed sign is $\eta_v=+1$.
There can be zero, one, or two inserted spheres; each loop is treated
independently. After undoing \eqref{eq:transport}, recovery returns the
physical block with the original vertex labels.

\subsection{Mercedes: one external insertion}

The genus-three Mercedes channel has a triangular rim and three spokes
meeting at a central vertex. Its four vertices and six internal edges
give $g=6-4+1=3$. Take spokes $1,2,3$ to be NS and rim $4,5,6$ to be
Ramond, with slots
\begin{equation}
\begin{array}{c|c|c}
\text{vertex}&(\infty,1,0)\text{ edges}&\text{form}\\ \hline
\text{center}&(1,2,3)&\rho_a\\
v_1&(1,4,6)&\rho_{f_1}^{(\eta_{v_1})}\\
v_2&(2,5,4)&\rho_{f_2}^{(\eta_{v_2})}\\
v_3&(3,6,5)&\rho_{f_3}^{(\eta_{v_3})}
\end{array}
\label{eq:mercedes-vertices}
\end{equation}
Order edges by number, orient each spoke from the center outwards, and
orient the rim $v_1\to v_2\to v_3\to v_1$. Take even intrinsic primary
parities, as in the saved Mercedes tests. Summing the local parity
constraints gives $a=f_1+f_2+f_3$ modulo two. The physical correction
\eqref{eq:transport} is simply
\begin{equation}
 (-1)^{f_1(1+\epsilon_6)+f_2(1+\epsilon_4)+f_3(1+\epsilon_5)}.
 \label{eq:mercedes-specialization}
\end{equation}
Only the parity of the third slot at each outer sphere appears.
There is one Ramond cycle, $(4,5,6)$, for which $K(C)=0$. Its physical
sector identity and obstruction are
\begin{align}
 (\eta_4\eta_5\eta_6)\star\mathbf F_{\hat\Gamma}
   &=-i\eta_{v_1}\eta_{v_2}\eta_{v_3}\mathbf F_{\hat\Gamma},
 \nonumber\\
 (-1)^{f_1+f_2+f_3}\eta_{v_1}\eta_{v_2}\eta_{v_3}&=-1
                           \quad\text{for an obstructed loop}.
 \label{eq:mercedes-obstruction}
\end{align}
The edge spin lifts and the three-point-form signs remain distinct.
For the ordinary and inserted calculations the auxiliary constants are
\begin{equation}
 \left.\mathbf F_{\F,\hat\Gamma}\right|_{q=0}
       =1+i\eta_4\eta_5\eta_6,\qquad
 \left.\mathbf F_{\F,\hat\Gamma}[\Theta\psi]_{\mathrm{diag}}
                                   \right|_{q=0}
       =\frac{1-i\eta_4\eta_5\eta_6}{\sqrt2}.
 \label{eq:mercedes-ground}
\end{equation}
This displays the switch of sector using a single insertion.

If the obstruction holds, split edge 4 with its original ket end at
$v_1$ and bra end at $v_2$, and place one external $\Theta\psi$ on the
new sphere. All three NS spokes remain sewn to the central sphere.
The enlarged calculation is a genus-three one-point block; its convolution
recovers the original physical block without an external insertion.
There are no additional external NS legs at the three outer vertices.
If the loop is unobstructed, use the original four spheres without the
insertion. These are the two Mercedes constructions tested below.

\section{External states and the full correlator}
\label{sec:external}

\subsection{Projecting an external state onto the branches}

An external tensor-PBW state need not be a branch primary. Expand it in
the same orthogonal double-Virasoro basis used on the internal edges.
At a Ramond puncture, its physical and auxiliary parities are
$\epsilon=A+|\alpha|$ and $\epsilon'=\mathsf A+\mathsf a$ modulo two.
Only branch primaries of total parity $\epsilon+\epsilon'$ can occur.
Using $G^{BB'}_h$ for the inverse Virasoro Gram matrix at weight $h$,
the BPZ expansion is
\begin{align}
 &\Psi_{-\mathsf A}u^{\mathsf a}\otimes\mathbb L_{-A}w^\alpha
 \nonumber\\
 &\quad=\sum_{n,B,C,B',C'}
 (G_{h_n^{(1)}})^{BB'}(G_{h_n^{(2)}})^{CC'}
 \Big\langle\!\Big\langle
 \frac{L_{-B'}^{(1)}L_{-C'}^{(2)}v_n^{\epsilon+\epsilon'}}
      {\|v_n^{\epsilon+\epsilon'}\|}
 \Big|\Psi_{-\mathsf A}u^{\mathsf a}\otimes\mathbb L_{-A}w^\alpha
 \Big\rangle
 \nonumber\\
 &\qquad\times
 \frac{L_{-B}^{(1)}L_{-C}^{(2)}v_n^{\epsilon+\epsilon'}}
      {\|v_n^{\epsilon+\epsilon'}\|}.
 \label{eq:external-descendants}
\end{align}
The superscript $\epsilon+\epsilon'$ is reduced to zero or one.
The finite sum is restricted by
$2n^2-1/8+|B|+|C|=|A|+|\mathsf A|$ and
$|B|=|B'|$, $|C|=|C'|$. This is the usual basis-expansion formula:
the overlap supplies a covariant component and the inverse Gram matrices
convert it to a basis coefficient. It does not require fitting the state
to an unknown linear combination.

For NS, replace the physical and auxiliary Ramond grounds by $\phi$ and
$\mathbf1$, replace $\mathbb L$ by $\mathbf L$, and use $v_n$ with
starting level $2n^2$. Its parity $2n+p$ must equal $\epsilon+\epsilon'$.
A ground tensor therefore involves only $n=0$ in NS or $n=\pm1/4$ in R,
with no Virasoro descendants. Insert the expansion at each external
puncture in \eqref{eq:double-virasoro} or
\eqref{eq:split-double-virasoro}. Whenever $B$ or $C$ is nonempty, the
corresponding Virasoro factor has that external descendant, which is
reduced by its own Ward identities. The primary norm roots must match
those in the branching coefficients.

\subsection{Combining holomorphic and antiholomorphic blocks}

Tildes denote the independent antiholomorphic data. For each choice of
local form labels, let $C_v$ be the full three-point structure constant
with the three holomorphic operators ordered before the three
antiholomorphic operators. The normalized chiral forms and BPZ metrics
are those already specified. Write $\mathcal O_a$ for the external full
fields in their prescribed order; $a,b$ below index these punctures,
not the form label $a_v$. The correlator in the plumbing
trivialization is
\begin{align}
 \left\langle\prod_a\mathcal O_a\right\rangle
 =\sum_{\substack{\text{internal representations and}\\
               \text{holomorphic/antiholomorphic vertex labels}}}{}&
 (-1)^{\sum_{v<w}(\sum_i\widetilde\epsilon_{v,i})
                         (\sum_j\epsilon_{w,j})
                +\sum_{a<b}\widetilde\epsilon_a\epsilon_b}
 \left(\prod_v C_v\right)
 \nonumber\\
 &\times\left(\prod_e q_e^{h_e}\bar q_e^{\widetilde h_e}\right)
             \mathbf F_{\hat\Gamma}\widetilde{\mathbf F}_{\hat\Gamma}.
 \label{eq:correlator}
\end{align}
The first sign moves the antiholomorphic operators at a vertex
past the holomorphic operators at later vertices. The second restores
the prescribed ordering of the external full fields. Each vertex's total
parity is fixed by \eqref{eq:vertex-selection}, so the sign can be taken
outside the descendant sums. For an even full three-point tensor the
holomorphic and antiholomorphic total parities agree; forbidden label
choices have $C_v=0$.

One can verify the sign directly before doing the descendant sums.
Graded factorization of the inverse pairings and three-point forms gives
\begin{align}
 &K(\epsilon+\widetilde\epsilon)-K(\epsilon)-K(\widetilde\epsilon)
 +\sum_e\epsilon_e\widetilde\epsilon_e
 +\sum_v\sum_{i<j}\widetilde\epsilon_{v,i}\epsilon_{v,j}
 \nonumber\\
 &\qquad=\sum_{v<w}\Bigl(\sum_i\widetilde\epsilon_{v,i}\Bigr)
                    \Bigl(\sum_j\epsilon_{w,j}\Bigr)
          +\sum_{a<b}\widetilde\epsilon_a\epsilon_b\pmod2.
 \label{eq:correlator-sign}
\end{align}
Each internal edge appears twice with equal parities. Moving these even
pairs as units cancels their internal contributions and leaves the vertex
and external-operator exchanges on the right. This accounts for the sign
in \eqref{eq:correlator} without assuming Ramond forms are absolute squares.

Sums over a continuous spectrum become integrals. A nontrivial primary
multiplicity metric contributes its inverse to the representation sum;
the physical Ramond ground metric is already in $\mathbb G$. A choice
of surface metric also supplies the common conformal-anomaly factor.
The block construction determines the chiral sewing coefficients, while
the spectrum and the $C_v$ are additional data of the SCFT.

\section{Directed numerical validation}
\label{sec:validation}

This section records all directed tests supporting these notes, including
the complete total-level-six extension in
Sec.~\ref{sec:total-six}. The tests distinguish
physical SCA PBW sewing, enlarged PBW sewing,
double Virasoro assembly, direct fermion sewing, and physical recovery.
The double Virasoro numerators are constructed from branching coefficients
and Virasoro blocks, without using the physical PBW answers. Numerical
discrepancies are reported as $|x-y|/\max(1,|x|,|y|)$; constrained zero
components are included. A vanishing test reports the absolute residual.
All results below are from saved computations; documentation changes do
not entail rerunning the numerical suites.

\subsection{Complete total level six: all 96 channel cases}
\label{sec:total-six}

Every channel in the following table was tested through total level six:
\begin{equation}
 r_e\ge0,\qquad \sum_e r_e\le6,
 \qquad r_e\in\tfrac12\mathbb Z\ (\NS),\quad
 r_e\in\mathbb Z\ (\R).
 \label{eq:total-six-domain}
\end{equation}
Thus an individual edge can reach level 6. The momenta and vertex-label
choices are specified in the channel descriptions below. Both NS primary parities
are included for NS--R--R theta; the other closed-block tests have even
intrinsic primary parities. All formal parity components, including
constrained zeros, are retained. These complete-domain calculations all use
40-digit arithmetic. Every one of the 96 cases passes.

\begin{center}
\small
\begin{tabular}{@{}lrrrr@{}}
\toprule
Channel & Cases & Multidegrees/case & Recovery error & Time (s)\\
\midrule
Theta NS--R--R & 16 & 140 & $2.09\times10^{-25}$ & 188.57\\
Theta all-NS & 2 & 455 & $9.29\times10^{-29}$ & 1.27\\
Glasses R/R & 8 & 140 & $2.87\times10^{-25}$ & 27.17\\
Glasses NS/R & 4 & 252 & $1.34\times10^{-25}$ & 4.73\\
Glasses NS/NS & 2 & 455 & $1.27\times10^{-29}$ & 0.58\\
Mercedes & 64 & 3906 & $4.81\times10^{-24}$ & 636.82\\
\bottomrule
\end{tabular}
\end{center}
Recovery errors compare the extracted physical block with independent
physical SCA PBW sewing. The times include the independent references,
comparisons, and coefficient output; they are validation-suite wall times,
not production block benchmarks. Some channels ran concurrently.

This extension includes the full CCY residues and the nonconstant
Schottky vacuum. The graph implementation caches pole geometry, fusion
polynomials, global vertices, and repeated recursion states. The universal
vacuum is evaluated from primitive Schottky cycles of the ordered graph.
NS--R--R vertices reuse the existing recursive branching tables, and
inserted vertices reuse the middle branching coefficients. The all-NS
vertices additionally require $n=\pm1,\pm\frac32$: their primary vectors
are constructed in free fields, converted to physical-PBW/auxiliary
coordinates, and evaluated by the local three-point form. These local
coefficients are cached. The physical reference uses full inverse Gram
matrices; in particular, all three Mercedes spoke metrics act on the
central tensor, including their off-diagonal entries.

For glasses and Mercedes, full split powers are retained in the domain
\begin{equation}
 \sum_{e:\text{unsplit}}r_e+
 \sum_{e:\text{split}}\max(r_{e_L},r_{e_R})\le6.
 \label{eq:total-six-split-domain}
\end{equation}
Its diagonal is exactly \eqref{eq:total-six-domain}, and it also contains
every expanded-graph multidegree whose ordinary total level is at most 6.
The unequal-power comparisons include nonzero modes of $\psi(1)$.
Theta retains its established split domain
$2r_1+r_L+r_R+2r_3\le12$ for
$q_1^{r_1}q_L^{r_L}q_R^{r_R}q_3^{r_3}$, with half-integer NS powers
and integer Ramond powers.

All 136 NS--R--R theta comparisons pass tolerance $10^{-20}$. Each
diagonal comparison covers 17920 components. The recovered-block error
is tabulated above; the restored double Virasoro numerator agrees with
direct enlarged PBW within $3.24\times10^{-25}$. Direct enlarged recovery
agrees with physical PBW within $6.83\times10^{-38}$, and the graph
convolution agrees with enlarged PBW within $9.80\times10^{-39}$.
All spin-lift evaluations agree within $2.50\times10^{-25}$ and the
physical loop-sector identity within $1.51\times10^{-37}$.
The 66304 full-split components agree with directly sewn split PBW within
$4.02\times10^{-22}$ for double Virasoro and $1.30\times10^{-38}$ for
the explicit convolution. The split fermion and enlarged diagonals agree
exactly with their zero-mode counterparts (7616 and 8960 components,
respectively). The 8960 uninserted obstructed components vanish within
$1.37\times10^{-37}$.

The other graph suites pass tolerance $10^{-18}$. Their forward
convolution, spin-lift, and independent split checks are:
\begin{center}
\small
\begin{tabular}{@{}lrrr@{}}
\toprule
Channel & Diagonal convolution & All spin lifts & Unequal split powers\\
\midrule
Theta all-NS & $4.34\times10^{-29}$ & $9.29\times10^{-29}$ & ---\\
Glasses R/R & $1.83\times10^{-24}$ & $5.73\times10^{-25}$ & $9.84\times10^{-24}$\\
Glasses NS/R & $8.75\times10^{-26}$ & $1.77\times10^{-25}$ & $4.35\times10^{-24}$\\
Glasses NS/NS & $3.93\times10^{-30}$ & $1.27\times10^{-29}$ & ---\\
Mercedes & $1.44\times10^{-23}$ & $6.79\times10^{-24}$ & $2.66\times10^{-23}$\\
\bottomrule
\end{tabular}
\end{center}
Each diagonal or spin-lift comparison covers 7280, 8960, 8064, 7280,
and 15998976 components, respectively. Unequal split powers add 36736
glasses R/R, 10752 glasses NS/R, and 11845632 Mercedes components.
In Mercedes each inserted case has 9690 full-split multidegrees, of which
5784 are off the diagonal. The 32 ordinary and 32 inserted Mercedes cases
both have maximum recovery error below $4.81\times10^{-24}$.
Uninserted obstructed numerators vanish within $1.40\times10^{-36}$,
$3.90\times10^{-37}$, and $1.96\times10^{-36}$ for glasses R/R,
glasses NS/R, and Mercedes, respectively.

A separate audit of the saved physical coefficients checks both the full
multidegree sets and the physical loop-sector identity, without recomputing
any block. It finds no missing or duplicated diagonal multidegree and
verifies the independent split-domain counts. The loop identities agree
within $1.78\times10^{-36}$ for glasses R/R, $1.47\times10^{-36}$ for
glasses NS/R, and $1.06\times10^{-38}$ for Mercedes. To verify the new
40-digit NS Ward implementation against its existing independent Python
implementation, all 910 all-NS theta case/multidegree pairs are also
compared with \texttt{DirectThetaOracle}. This separate machine-precision
reference agrees within $5.78\times10^{-12}$, at tolerance $10^{-9}$;
it does not determine the errors in the 40-digit tables above.

Before the complete runs, the extended implementation also passed
total-level-two checks for all-NS theta (two cases, 35 multidegrees each),
glasses R/R (four cases, 14 multidegrees each), and Mercedes (four cases,
71 multidegrees each). Their saved results are retained with the extension.

The extension exposed a sign error in the new graph CCY driver that the
initial cutoff could not detect. For a self-loop sewn between local slots
$(1,0)$, the residue at null level $rs$ needs the factor $(-1)^{rs}$ from
the middle-slot Virasoro Ward normalization. Omitting it first produces
a discrepancy at physical multidegree $(0,0,3)$ in glasses. All glasses
results above are from fresh runs after correcting that factor. The
Mercedes graph has no self-loops, so its concurrent run does not execute
the corrected branch. No production header or manuscript formula was
changed by this test-driver correction.

The Mercedes run sums 29240 branching tuples and makes 896784 CCY
recursion transitions. Its double Virasoro stage took 371.45 seconds,
direct auxiliary sewing 0.20 seconds, and physical PBW sewing 29.80 seconds
across all cases. The remaining time is coefficient comparisons, recovery,
spin evaluations, and output. It saves all 249984 diagonal coefficient
rows, with physical, enlarged, and recovered values.

All extension sources, results, coefficients, and reproduction commands
are in \path{C++/experiments/total_level6_2026-09-16/}.
The \path{README.md} gives the cutoffs and implementation details;
\path{manifest.json} summarizes the results and source hashes.
The script \path{run_all.py} performs a fresh sequential reproduction.
The saved-data audit is \path{coverage_and_sector_results.json};
the separate NS Ward-port check is \path{ns_physical_port_results.json}.
The diagnostic self-loop failures are retained separately and are not
counted as passing results. These tests use the NS/R assignments
and primary-parity cases listed below, at the same generic momenta; they do
not add the inequivalent four-edge Ramond cycle on the Mercedes graph or singular limits.



\subsection{Initial NS--R--R theta tests, including independent split powers}

The NS--R--R theta tests use 40-digit arithmetic. We take $b=7/5$ and
graph momenta $(11/23,13/29,17/31)$. Both $p_1,f=0,1$, all four vertex-sign
pairs, every formal parity component, and all eight evaluations of the
spin lifts are covered, giving 16 cases. The diagonal physical cutoff is
total level 3: 30 multidegrees and 240 parity components per case.
The independent split series retains
$q_1^{r_1}q_L^{r_L}q_R^{r_R}q_3^{r_3}$ with
$2r_1+r_L+r_R+2r_3\le6$, half-integer NS powers, and integer Ramond
powers; its diagonal is the same physical cutoff. There are 130 multidegrees and
1040 components per inserted case. The full split numerator is
checked against an independently sewn enlarged PBW numerator, including
the modes of $\psi(1)$ that change the middle level. The double Virasoro
calculation uses both CCY recursions and the Schottky vacuum. The numerical
insertion is $Q\Theta\psi$ in both the numerator and auxiliary factor.

{\small
\begin{longtable}{@{}p{0.60\linewidth}rr@{}}
\toprule
Comparison & Components & Maximum discrepancy\\
\midrule
\endhead
Recovered double Virasoro vs physical PBW & 3840 & $7.32\times10^{-28}$\\
Double Virasoro numerator vs enlarged PBW & 3840 & $2.20\times10^{-27}$\\
Recovery from enlarged PBW vs physical PBW & 3840 & $1.01\times10^{-40}$\\
Graph convolution vs enlarged PBW & 3840 & $7.73\times10^{-41}$\\
Recovered block vs PBW at all spin lifts & 3840 & $1.35\times10^{-27}$\\
Physical loop-sector identity & 3840 & $1.04\times10^{-40}$\\
Split fermion diagonal vs zero-mode factor & 1536 & $0$\\
Split enlarged diagonal vs zero-mode enlarged PBW & 1920 & $0$\\
Full split double Virasoro vs split enlarged PBW & 8320 & $7.06\times10^{-26}$\\
Full split convolution vs split enlarged PBW & 8320 & $6.29\times10^{-41}$\\
Uninserted enlarged PBW in obstructed cases vs zero & 1920 & $4.34\times10^{-40}$\\
\bottomrule
\end{longtable}}
All 136 comparisons pass tolerance $10^{-20}$. The two diagonal/zero-mode
identities agree exactly in the saved arithmetic. The 16 fresh numerical
processes took 3.61 seconds together; this is a validation-suite time,
not a high-level block benchmark.

\subsection{Initial all-NS theta and local all-NS vertices}

Exact rational tests cover all eight intrinsic-primary-parity assignments
and all eight branch-label tuples $n_i\in\{0,\frac12\}$. All 64 local
vertex comparisons and all 384 $L_{-1}$ Ward identities (three slots in
each Virasoro copy) pass. They use the all-NS enlarged-vertex convention
stated above, including its operator-moving signs.

The independent closed all-NS theta calculation has even primaries and
total level 2. All 35 coefficients agree with physical PBW within
$4.67\times10^{-15}$. All eight spin-lift choices in both vertex-parity
sectors give 16 evaluation checks, with maximum discrepancy
$5.56\times10^{-17}$. The tolerance is $10^{-10}$. The series backend
uses machine arithmetic, with 40-digit branching-polynomial evaluation;
this is not a 40-digit end-to-end calculation.

\subsection{Initial glasses tests with zero, one, or two insertions}

The glasses tests use the same momenta (bridge, left handle, right handle),
machine arithmetic, even intrinsic primary parities, and individual edge
level 1. They cover R/R, NS/R,
and NS/NS handles; R/NS is the isomorphic assignment obtained by exchanging
the named handles, rather than a separate numerical run.
Both vertex-form parities, all Ramond signs, and every spin-parity slot
are included: eight R/R cases, four NS/R cases, and two NS/NS cases.
The R/R tests include zero, one, and two inserted loops;
960 of the full-split components have unequal powers on a split pair.
At this cutoff every Virasoro edge has descendant level at most one, so
the exact CCY answer equals its global term. These glasses tests do not
test higher-level residues or nonconstant Schottky corrections.

\begin{center}
\small
\begin{tabular}{@{}p{0.60\linewidth}rr@{}}
\toprule
Comparison & Components & Maximum discrepancy\\
\midrule
R/R diagonal numerator vs PBW convolution & 768 & $7.02\times10^{-13}$\\
R/R recovery vs physical PBW & 768 & $8.71\times10^{-14}$\\
R/R full split numerator vs PBW convolution & 1728 & $8.30\times10^{-13}$\\
NS/R numerator vs PBW convolution & 576 & $3.17\times10^{-14}$\\
NS/R recovery vs physical PBW & 576 & $2.79\times10^{-14}$\\
NS/NS numerator vs PBW convolution & 432 & $1.45\times10^{-15}$\\
NS/NS recovery vs physical PBW & 432 & $1.90\times10^{-15}$\\
\bottomrule
\end{tabular}
\end{center}
The uninserted enlarged block is separately checked to vanish in every
obstructed R/R case, with maximum absolute residual
$1.17\times10^{-13}$. All checks pass tolerance $10^{-9}$.
The R/R diagonal, full-split, and NS-assignment tests record runtimes of
0.079, 0.105, and 0.042 seconds, respectively, excluding compilation.

\subsection{Initial Mercedes tests: ordinary and obstructed sectors}

For the graph in \eqref{eq:mercedes-vertices}, take $b=7/5$ and edge
momenta $(11/23,13/29,17/31,19/37,23/41,29/43)$. Every original edge has
an independent physical cutoff of level 1. The three NS powers are
$0,\frac12,1$, and the three Ramond powers are $0,1$, giving
$3^3 2^3=216$ multidegrees per case, including total level 6.
The inserted test also retains unequal split powers $(0,1)$ and $(1,0)$,
with each segment independently bounded by level 1. Thus each inserted
case has 432 full-split multidegrees, half of them diagonal.

All eight choices of $(f_1,f_2,f_3)$ and all eight choices of
$(\eta_{v_1},\eta_{v_2},\eta_{v_3})$ are tested. Equation
\eqref{eq:mercedes-obstruction} gives 32 ordinary and 32 obstructed cases.
Every case retains all 64 formal parity components and compares all 64
evaluations of the edge spin lifts. Intrinsic primary parities are even.
Arithmetic is complex machine precision; the central NS branching
polynomials are evaluated at 40 digits before conversion.

The physical reference contracts SCA PBW vertices and inverse SCA Gram
matrices. Its Ramond vertices use \texttt{ScaWard}, and its central NS
vertex uses the independent \texttt{NSDescendantThreeForm}; it does not
use branching data. The enlarged numerator instead uses
\texttt{LowAnchors} at the three NS--R--R vertices,
\texttt{ns\_fusion\_data} at the center, and \texttt{MiddleBranching} on
the split edge, followed by two Virasoro graph contractions. Local vertex
and branching values are reused within the run. At this cutoff every
Virasoro descendant edge has level at most one, so its exact CCY answer
is the global term, evaluated with the $L_{-1}$ Ward formulas.
No higher Kac residue or nonconstant Schottky correction enters this test.

The auxiliary factor is sewn directly from fermion Fock states and Ward
forms. Recovery uses only the double Virasoro numerator and this factor:
it subtracts lower-level convolutions, divides by the matching-sector
constant, and undoes the parity correction \eqref{eq:transport}. Physical PBW coefficients
are used only for comparison. Unequal split powers are compared with the
independent physical-PBW/fermion convolution, including nonzero fermion
modes. Both inserted factors use $Q\Theta\psi$, so the normalization
cancels upon recovery.

\begin{center}
\small
\begin{tabular}{@{}p{0.53\linewidth}rr@{}}
\toprule
Comparison & Ordinary & Inserted\\
\midrule
Recovered physical block vs PBW & $1.55\times10^{-12}$ & $1.60\times10^{-12}$\\
Diagonal numerator vs PBW convolution & $1.26\times10^{-12}$ & $2.31\times10^{-12}$\\
Physical loop-sector identity & $6.67\times10^{-16}$ & $6.67\times10^{-16}$\\
Recovered block vs PBW at all spin lifts & $2.14\times10^{-12}$ & $2.20\times10^{-12}$\\
Unequal split powers vs PBW convolution & --- & $2.03\times10^{-12}$\\
\bottomrule
\end{tabular}
\end{center}
Each of the first four rows covers 442368 components per column:
$32\times216\times64$. The recovery comparisons contain 6912 nonzero
formal components per column; constrained zeros are retained.
The unequal-split check adds 442368 components. The uninserted enlarged
blocks in all 32 obstructed cases vanish with maximum absolute residual
$9.14\times10^{-13}$, while their physical blocks are generally nonzero.
All checks pass tolerance $10^{-8}$.

The new central-fermion adapter initially omitted the bra-slot BPZ sign.
The raw sphere Pfaffian needs the factor $(-1)^{\epsilon^{\prime}_1}$, consistently
with the auxiliary $\psi_{-1/2}$ norm $-1$ used above. Correcting that
adapter resolved the mismatch; the graph formulas and production headers
were unchanged. The results tabulated here use the corrected convention.

The C++ test process took 1.93 seconds of wall time, including coefficient
output; its internal timer recorded 1.59 seconds. Central-data export took
0.098 seconds and compilation 3.29 seconds, for a fresh total of
5.32 seconds. These are low-level validation timings. The run saves all
13824 diagonal coefficient rows, each with the physical PBW, enlarged
double Virasoro, and recovered parity arrays.

\subsection{Local Ramond identity and external-state projections}

The Ramond odd-ground identity used in \eqref{eq:physical-sector} passes
864 local Ward tests at 40 digits. These cover $p,f=0,1$, both vertex
signs, the NS primary and its $G_{-1/2}$ and $L_{-1}$ descendants, and
both Ramond grounds and their $G_{-1}$ and $L_{-1}$ descendants in each
Ramond slot. The maximum discrepancy is $2.90\times10^{-41}$.

At $b=7/5$, $P=13/29$, 40-digit external-state BPZ projections reconstruct
every tensor-PBW input at NS levels $0,\frac12,1$ and R levels $0,1$:
\begin{center}
\small
\begin{tabular}{@{}lrrrr@{}}
\toprule
Sector, level & Inputs & Gram entries & Gram error & Reconstruction error\\
\midrule
NS, $0$ & 1 & 1 & $0$ & $0$\\
NS, $\frac12$ & 2 & 4 & $1.40\times10^{-32}$ & $5.68\times10^{-33}$\\
NS, $1$ & 2 & 4 & $1.42\times10^{-32}$ & $3.62\times10^{-33}$\\
R, $0$ & 4 & 16 & $5.86\times10^{-33}$ & $2.93\times10^{-33}$\\
R, $1$ & 12 & 144 & $1.27\times10^{-30}$ & $1.29\times10^{-30}$\\
\bottomrule
\end{tabular}
\end{center}
This gives 21 independent inputs and 169 branch-Gram entries, including
all four Ramond ground tensors. Branch vectors are obtained from
free-field primary and generator actions. Their coordinate conversion to
the physical-PBW/auxiliary basis has maximum residual
$2.70\times10^{-32}$ and is used only to evaluate BPZ overlaps.
The final reconstruction uses independently known inverse Virasoro Gram
matrices, not a fit of the input to the branch basis. This checks the
external-state expansion locally, not a full correlator with external
insertions.

\subsection{Exact graph-sign and parity-algebra checks}

Exact operator-order permutations verify \eqref{eq:correlator-sign}, while
independent parity enumeration checks the associativity of the graph
convolution:
\begin{center}
\small
\begin{tabular}{@{}p{0.52\linewidth}rr@{}}
\toprule
Graph & Parity pairs & Associativity triples\\
\midrule
Theta & 64 & 512\\
Glasses & 64 & 512\\
Local ordering with external states & 4096 & 262144\\
\bottomrule
\end{tabular}
\end{center}
The theta driver also checks all eight half-edge parity assignments and
all 64 explicit convolution-sign pairs. Four loop cases check commutation of the cycle monomials
and the squares of the normalized cycle terms in \eqref{eq:constant-projectors}, and 64 local NS--R--R cases check phase
factorization. All finite checks pass exactly. The external-state row is a local
ordering check before sewing the NS slots; it is not a separate recovery
construction. The Mercedes calculation retains its central sphere and
uses the single external insertion described above.

\subsection{All-NS Mercedes: comparison with NS \texorpdfstring{$c$}{c}-recursion}

We also compare the all-NS Mercedes block with the NS $c$-recursion in
the SCblock notes, independently of physical PBW sewing. All six edges
are NS, with even intrinsic primary parities. Use the vertex orders in
\eqref{eq:mercedes-vertices}, $b=7/5$, and momenta
$(11/23,13/29,17/31,19/37,23/41,29/43)$. Both calculations use 40-digit
arithmetic and retain every nonnegative half-integer multidegree with
total level at most 6: 18564 multidegrees. All eight allowed vertex-form
assignments are included. At fixed NS multidegree the edge parities are
fixed, so coefficientwise agreement also covers all 64 spin-lift choices.

The double-Virasoro side uses both full CCY recursions and Schottky vacuum
factors, then removes the directly sewn auxiliary fermion factor. The
independent NS recursion uses the Kac poles, inverse null slopes, fusion
polynomials, and explicit global $\mathfrak{osp}(1|2)$ vertices. Its vacuum
seed is expanded from the primitive Schottky cycles of this graph. Through
level 6, the four triangles contribute supercurrent terms at levels
$9/2,11/2$ and stress-tensor terms at level 6; the three four-edge cycles
first contribute at level 6. Thus this comparison includes the nonconstant
vacuum seed.

The vacuum/global product requires the full sign in \eqref{eq:star}, with
$\epsilon$ the global parity and $\epsilon'$ the vacuum parity. In
particular, moving the middle-slot odd global operators past the ket-slot
odd vacuum operators contributes
$(-1)^{\sum_v\epsilon_{v,2}\epsilon'_{v,3}}$ in addition to the graph
permutation sign. This term cancels in theta but survives in Mercedes.
Omitting it gives 130 discrepant coefficients, beginning at total level 5;
including it gives agreement throughout the requested domain. The maximum
scaled error is $1.05\times10^{-28}$ and the maximum absolute error is
$1.18\times10^{-28}$.

The fresh run takes 15.78 seconds: 10.03 seconds for the enlarged
double-Virasoro blocks, 0.013 seconds for the auxiliary factor, 0.162 seconds
for recovery, and 5.58 seconds for NS recursion, comparisons, and output.
Compilation is excluded. Sources, coefficients, the diagnostic with the
omitted sign, and reproduction instructions are saved in
\path{C++/experiments/mercedes_all_ns_c_recursion_2026-09-16/}.
This additional test covers the all-NS assignment with even intrinsic
primaries; it does not extend the tested mixed NS/R assignments.

\paragraph{Extension to total level 10.}
At the same parameters and precision, the two methods also agree on all
230230 multidegrees through total level 10. The all-even vertex-form sector
contains 29029 multidegrees, and each of the other seven sectors contains
28743. The maximum scaled error is $2.05\times10^{-25}$ and the maximum
absolute error is $4.00\times10^{-25}$; no coefficient fails the
$10^{-18}$ tolerance. No additional physical PBW block is computed.

This extension uses the full truncated Schottky product: 13 primitive
unoriented cycles of lengths at most six and 20 fermionic oscillator
factors, including their products. The lift phases are transported under
cycle composition, and each half-multiplier is checked to square to the
corresponding multiplier in the cycle convolution algebra. The chosen
coordinate matrices give edge-map determinant $+q_e$ when exactly one
endpoint is an infinity slot, and $-q_e$ otherwise. This replaces the old
generic helper's uniform $-q_e$ assumption; its first effect on the
bosonic vacuum is at total level 9, so the level-six result is unaffected.

The fresh calculation retains 45405 branch tuples and performs 10921308
CCY residue transitions. It generates the residual Virasoro level domain
directly after subtracting each branch's starting levels, and convolves
only pairs within that domain. Branching data and local recursion data
are reused. Recovery updates only the higher coefficients that receive
a nonconstant auxiliary contribution. On the NS recursion side, the
full-block cache is cleared between target multidegrees: accumulated
shifts plus remaining levels equal the target, so those entries cannot
be shared by different targets. The reusable local caches are retained.

\begin{center}
\small
\begin{tabular}{@{}lr@{}}
\toprule
Stage, 40-digit arithmetic & Time (s)\\
\midrule
Double-Virasoro enlarged blocks & 681.135\\
Direct auxiliary factor & 0.183\\
Collect enlarged coefficients & 0.071\\
Recovery and coefficient output & 1.892\\
\midrule
Full double-Virasoro pipeline & 683.282\\
NS Schottky seed & 0.178\\
NS recursion and coefficient output & 270.148\\
Full NS recursion stage & 270.326\\
\bottomrule
\end{tabular}
\end{center}
The two sequential computational stages take 953.61 seconds in total,
excluding compilation and the subsequent comparison of saved coefficients.
Both coefficient sets, errors at each total half-level, source hashes, and
reproduction instructions are saved in
\path{C++/experiments/mercedes_all_ns_level10_2026-09-16/}.

\subsection{Saved evidence, reproduction, and scope}

The theta, glasses, local Ward, external-state, and finite graph scripts,
commands, source hashes, and results are saved in
\path{C++/experiments/arbitrary_plumbing_2026-09-16/}.
The reproduction guides are \path{theta_README.md} and
\path{glasses_README.md}; the directory's \path{README.md} maps the formulas
to their checks. The result files are \path{theta_summary.json},
\path{theta_all_ns_results.json}, \path{theta_ramond_local_results.json},
\path{external_state_results.json}, \path{glasses_results.json},
\path{glasses_full_split_results.json},
\path{glasses_ns_assignments_results.json}, and
\path{graph_sign_results.json}.

The Mercedes sources and results are saved separately in
\path{C++/experiments/genus3_mercedes_2026-09-16/}.
Running \texttt{python3 run.py} from that directory exports the central
data, compiles the C++ driver, and executes the complete 64-case suite.
The \path{manifest.json} records parameters, graph ordering, source and
result hashes, timings, and aggregate errors. The \path{results.json}
records every case and its worst discrepancies; \path{coefficients.jsonl}
contains the diagonal coefficient arrays, and \path{central_manifest.json}
records the independent central-vertex data. Its \path{README.md} explains
the conventions and reproduction commands.

These checks concern generic momenta at the stated cutoffs. For Mercedes
they cover the all-NS and triangular Ramond assignments, excluding the
inequivalent four-edge Ramond cycle and intrinsic-odd NS primaries. They do not test singular
weights, full external-Ramond correlators, or a SCFT spectral sum or moduli
integral. No arbitrary-genus numerical frontend is claimed: the general
argument uses the local identities and graph contractions above.


\end{document}
